% sections/ndss_02_background.tex  — NDSS 2027

\subsection{LLM Agent Skill Marketplaces}

LLM skill marketplaces expose installable tools that extend agent
capabilities via standardized protocols such as the Model Context Protocol (MCP)~\cite{anthropic2024mcp}.
A skill is a self-contained package comprising:
(i) a \emph{description} (natural-language documentation visible to users),
(ii) \emph{tool definitions} (structured capability declarations),
and (iii) \emph{implementation code} (executed by the agent runtime).
\clawhub, launched with MCP in late 2024, is the largest such marketplace,
with \NskillsCH skills and $>$2.1M installs as of March 2026.

\paragraph{Ambient authority and sessions.}
Installed skills execute with the privileges of the agent runtime---typically
the user's OS account.
Skills can read and write files, execute shell commands, make network requests,
and read credentials without per-action OS-level confirmation.
This \emph{ambient authority} model, common in agent frameworks~\cite{ruan2024identifying},
means that a skill's behavior is constrained only by its own code
and the LLM's willingness to call its tools.

\subsection{Threat Model}

\paragraph{Principals.}
We consider three principals:
(1)~the \emph{user}, who installs and interacts with skills;
(2)~the \emph{skill developer}, who publishes skills to the marketplace;
and (3)~the \emph{attacker}, who may be a malicious skill developer
or an attacker controlling content fetched by a legitimate source skill.

\paragraph{Attacker capabilities.}
The attacker:
\begin{itemize}[leftmargin=*,topsep=2pt]
  \item can publish skills with benign-looking descriptions
    that omit dangerous capabilities (supply-chain attacker), OR
  \item can control content deliverable through a source skill's
    data-fetch operations (\eg web pages, emails, API responses)---the
    \emph{prompt injection via external content} model~\cite{greshake2023indirect}.
\end{itemize}

\paragraph{Attack goal.}
The attacker aims to cause the agent to invoke a dangerous capability
$\kappa$ of an installed sink skill, without the user's knowledge
or direct authorization, producing a real system-level effect $e$
(\eg file deletion, shell command execution, data exfiltration).

\paragraph{Out of scope.}
We do not model: direct UI-level social engineering of users,
OS-level privilege escalation beyond the agent runtime's account,
or attacks on the marketplace platform infrastructure itself.

\paragraph{Key assumption.}
The attack is \emph{enabled by the IC-1 gap}:
the user installed the sink skill because its description appeared benign,
and therefore has no awareness that $\kappa$ is possible.
This distinguishes our threat from known prompt injection scenarios
where the dangerous capability is explicitly documented.

\subsection{Why Descriptions Are the Only Defense}

Unlike Android (permission manifests~\cite{felt2011android}),
npm (lockfiles + audit~\cite{synk2023report}),
or sandboxed browser extensions (capability APIs),
LLM skill marketplaces have \emph{no machine-enforceable permission model}.
The MCP standard defines tool schemas, but these describe the API surface,
not the security implications of tool calls.
A skill declaring a \texttt{run\_command(cmd: str)} tool provides no
machine-readable signal that \texttt{cmd} is user-controlled or dangerous.
Natural-language descriptions are therefore the \emph{de facto}
capability disclosure mechanism, and their accuracy is a security property.
